% Two domains of the same circle parametrization; transitions have two components.
\begin{tikzpicture}[figure picture,x=1mm,y=1mm,
  every node/.append style={fill=none,inner sep=1pt}]
 \path[use as bounding box] (0,0) rectangle (112,83);
 \node at (28,78) {$\chi_1:\ \theta\in(-\pi,\pi)$};
 \node at (84,78) {$\chi_2:\ \theta'\in(0,2\pi)$};
 \foreach \cx/\circlecolor in {28/FigureBlue,84/FigureGreen}{
   \draw[\circlecolor,line width=.85pt] (\cx,53) circle[radius=16mm];
   \fill[FigureInk] (\cx,69) circle[radius=.7mm];
   \fill[FigureInk] (\cx,37) circle[radius=.7mm];
   \node[above right] at (\cx,69) {$p_+$};
   \node[below right] at (\cx,37) {$p_-$};
 }
 \filldraw[fill=white,draw=FigureBlue,line width=.85pt] (12,53) circle[radius=1mm];
 \node[below right] at (12,52) {$(-1,0)$};
 \filldraw[fill=white,draw=FigureGreen,line width=.85pt] (100,53) circle[radius=1mm];
 \node[below left] at (100,52) {$(1,0)$};
 \node at (28,27) {$p_+:\ \theta=\pi/2$};
 \node at (84,27) {$p_+:\ \theta'=\pi/2$};
 \node at (28,16) {$p_-:\ \theta=-\pi/2$};
 \node at (84,16) {$p_-:\ \theta'=3\pi/2$};
 \draw[FigureMuted,line width=.6pt,-{Stealth[length=1.8mm]}] (48,27)--(64,27);
 \draw[FigureMuted,line width=.6pt,-{Stealth[length=1.8mm]}] (48,16)--(64,16);
 \node[figure note] at (56,3) {$\theta>0:\ \theta'=\theta;\qquad\theta<0:\ \theta'=\theta+2\pi$};
\end{tikzpicture}
